\section{Постановка задачи}
\label{sec:task}

Целью лабораторной работы является ознакомление, анализ и приобретение навыков реализации модели релаксационной сети для задачи распознавания или классификации образов.

Реализуется универсальная модель сети Хопфилда, покрывающая все комбинации типа состояния и режима дискретного времени из вариантов задания: дискретное и непрерывное состояние, синхронный и асинхронный режимы обновления. В отчёте приведены результаты для всех четырёх режимов.

Требуется:
\begin{enumerate}
	\item построить модель релаксационной сети, решающую задачу распознавания образов;
	\item построить семейства графиков зависимостей параметров модели;
	\item дать ответы на вопросы из списка;
	\item предоставить письменный отчёт.
\end{enumerate}

Согласно требованиям методических указаний:
\begin{itemize}
	\item для сети Хопфилда с непрерывным состоянием используется гиперболический тангенс в качестве функции активации;
	\item для сети Хопфилда с дискретным состоянием используется модифицированная функция знака;
	\item веса рассчитываются методом проекции или методом дельта-проекции;
	\item предусмотрена проверка достижения состояния релаксации, не допускающая зацикливания алгоритма, и подчёт числа итераций до релаксации;
	\item распознаваемые образы кодируются значениями $1$ или $-1$ (в зашумлённых образах допускается также значение $0$).
\end{itemize}

Эталонные образы~\mdash  растровые битмапы букв и цифр размера $10 \times 10$, преобразованные в векторы длины $n = 100$ с кодированием
\begin{equation}
	x_i =
	\begin{cases}
		1, & \text{пиксель «включён»}, \\
		-1, & \text{пиксель «выключен»}.
	\end{cases}
\end{equation}
Зашумление реализуется знакопереворотом: каждая компонента с вероятностью $p$ умножается на $-1$.
